\section*{Resumen}\addcontentsline{toc}{section}{Resumen}
El grupo desplegó desde cero una celda 5G NR Standalone (opción~2, Release~15) con una USRP B210, el gNB OCUDU (ex srsRAN Project) y el núcleo Open5GS, todo dentro de contenedores Docker y en una jaula de Faraday. La celda principal opera en la banda n78 (TDD, 3525~MHz, 20~MHz, $\mu=1$, 51~PRB) con una frecuencia de muestreo de 23,04~MS/s, es decir, una FFT de 768~puntos a 30~kHz (equivalente a $1536\times 15$~kHz). También se puso en marcha una celda n3 FDD (1842,5/1747,5~MHz, 15~kHz, 106~PRB).

Un celular comercial (Xiaomi Redmi Note~14~5G) con una SIM de laboratorio programada por el grupo se registró en el núcleo, obtuvo una dirección IP y navegó por Internet a través de NAT. Para llegar a una conexión estable hubo que resolver cuatro problemas: el PLMN (con 722/999 los teléfonos no generaban PRACH; con el PLMN de prueba 001/01 se registraron al primer intento), el CSI que el teléfono no reportaba, una corrección de reloj que desplazaba la portadora sin corregir el muestreo y provocaba una deriva continua del \emph{timing advance}, y un corte del lado del teléfono cada 60~s, que se mitigó con el modo ``NR only'' y cuya relación con el anuncio de voz IMS sigue siendo una hipótesis.

Con la celda recién arrancada, la bajada TCP alcanzó $\approx 54$~Mbit/s (tasa del emisor), un 75~\% de la referencia aproximada de 71,7~Mbit/s que se obtiene de TS~38.306 para el patrón TDD usado. La subida llegó a 14,8~Mbit/s, contra una referencia de $\approx 35$~Mbit/s. El RTT mínimo fue de 16,7~ms y el mediano de 27,7~ms, aunque esta medición se hizo con la celda ya degradada. DFT-s-OFDM funcionó en la subida, pero no mejoró el rendimiento en este escenario de SINR alta. Eso es coherente con la teoría si el UE no estaba limitado en potencia, algo que queda como hipótesis porque no se registró el PHR. El informe desarrolla la teoría de capa física y de métodos de acceso que explica cada una de estas observaciones.
